\documentclass[10pt,a4paper]{amsart}
\usepackage[margin=19mm]{geometry}
\usepackage{amsmath,amssymb,mathtools}
\usepackage[scr=rsfs,cal=euler]{mathalfa}
\usepackage{xfrac}
\usepackage[unicode,hidelinks,bookmarks=false]{hyperref}
\newcommand{\ie}{i.e.\ }
\newcommand{\cf}{cf.\ }
\newcommand{\resp}{respectively\ }
\newcommand{\Subsection}{\subsection}
\providecommand{\nobreakdash}{\nobreak\hbox{-}}
\setlength{\emergencystretch}{2em}
\multlinegap0pt
\usepackage{bm}
\usepackage{bbm}
\usepackage{dsfont}
\usepackage{xspace}
\usepackage{cancel}
\usepackage{mathtools}
\usepackage{amsthm}
\makeatletter
\@ifundefined{theorem}{\newtheorem{theorem}{Theorem}}{}
\@ifundefined{lemma}{\newtheorem{lemma}[theorem]{Lemma}}{}
\makeatother
\makeatletter
\expandafter\ifx\csname claimproof\endcsname\relax
\newenvironment{claimproof}[1][\proofname]{\proof[#1]\renewcommand{\qedsymbol}{$\blacksquare$}}{\endproof}
\fi
\renewcommand{\phi}{\ensuremath{\varphi}}
\expandafter\ifx\csname subclaimproof\endcsname\relax
\newenvironment{subclaimproof}[1][\proofname]{\proof[#1]\renewcommand{\qedsymbol}{$\blacklozenge$}}{\endproof}
\fi
\makeatother
\title[Critical edge sets in vertex-critical graphs]{Critical edge sets in vertex-critical graphs}
\author{Ema Skottova, Raphael Steiner}
\date{Source: arXiv:2508.08703v1 (CC BY 4.0)}
\begin{document}
\maketitle

\noindent\emph{Source and licence.} Critical edge sets in vertex-critical graphs. Version arXiv:2508.08703v1 (CC BY 4.0), distributed under the Creative Commons Attribution 4.0 International licence, as recorded in the arXiv licence metadata for this exact version and shown on the abstract page https://arxiv.org/abs/2508.08703v1. Full rights notice: licenses/arxiv-2508.08703-CC-BY-4.0.txt.\par\smallskip

\noindent\textit{Editorial scope.} Counted unit: the Lemma whose source label is \texttt{colorability-lemma} in the article (source0.tex lines 382--384, source label \texttt{colorability-lemma}), delivered together with the article's own complete original proof of it (source0.tex lines 394--413, one \texttt{proof} environment of 20 lines). No mathematics is written, repaired or re-derived: statement and proof are the article's own text line for line. The proof is detached from the statement in the source: the article states the result at lines 382--384 and prints its proof at lines 394--413, and the proof environment's own header names the statement, so the association is the article's own. Required context retained uncounted: main-thm (lines 374--376); partial-periodicity-lemma (lines 388--390). Every definition, display and label of those lines is retained unchanged. Omitted from this package and not redistributed: 1--373, 377--381, 385--387, 391--393, 414--1455. Nothing inside a retained excerpt is omitted or altered: every retained body is delivered exactly as the article's lines are written, so no omission receipt of any kind accompanies this package. Citations: the retained excerpts carry no citation command at all, so no bibliography entry of the article is transcribed and no reference is redistributed. Reference adaptation: none was needed and none was made; every reference inside the delivered unit resolves inside it, so no unresolved reference or citation is produced. Printed slips kept literal: no printed word, symbol, hypothesis or number of the counted statement or of its proof was altered. Layout and class: the article's own class is \texttt{amsart} with packages including graphicx, geometry, hyperref, amsthm, todonotes, thm-restate, amssymb, amsmath, comment, xcolor, blkarray, lmodern, babel, float; the wrapper is the lane's pinned \texttt{amsart} editorial wrapper at 10pt on A4, which restates the theorem-like environments the retained text needs and adds the article's own macro definitions that the retained text needs (\texttt{\textbackslash phi}), with extra wrapper packages (bm, bbm, dsfont, xspace, cancel, mathtools). The delivered numbering is the unit's own. Assets: no retained span contains a figure environment, a graphics-inclusion command, a PDF-page inclusion, a table or a TikZ picture; no image file is copied into this package, the package declares no asset dependency and no image file is redistributed with it. Licence: version arXiv:2508.08703v1 of this article is distributed under the Creative Commons Attribution 4.0 International licence, as recorded in the arXiv licence metadata for this exact version and shown on the abstract page https://arxiv.org/abs/2508.08703v1. A full rights notice is delivered at licenses/arxiv-2508.08703-CC-BY-4.0.txt. Characters: every retained excerpt is pure ASCII, so the delivered engine, XeLaTeX with its default Latin Modern text font, reports no missing character.
    \begin{theorem}\label{main-thm}
    Let $k , r, q$ and $m$ be integers such that $k\ge 5$, $r\ge 1$, $q \geq 24r + 12$ is a multiple of $4$ and $m > 18r+2$. Then $G_{k, m, q}$ is a $(k,r)$-graph.
    \end{theorem}

    \begin{lemma}[Partial Periodicity]\label{partial-periodicity-lemma}
        Let $k, r, m, q$ be as in Theorem~\ref{main-thm}. Let $G'=G_{k, m, q}-R$ be a graph obtained from $G_{k, m, q}$ by deleting an edge set $R$ with $|R|\le r$. Let $S$ be a set of at least $3n_{k, m}$ consecutive unaffected vertices in $G'$. Let $S'$ be a set of at most $\frac{N-1}{4}$ consecutive vertices containing $S$. Then every proper $(k - 1)$-coloring of $G'$ is periodic on $S'$.
    \end{lemma}

    \begin{lemma}[Colorability]\label{colorability-lemma}
        Let $k, r, m, q$ be as in Theorem~\ref{main-thm}. For every $v \in V(G_{k, m, q})$, the graph $G_{k, m, q} - v$ is $(k-1)$-colorable.
    \end{lemma}

    \begin{proof}[Proof of Theorem~\ref{main-thm}, assuming Lemma~\ref{colorability-lemma} and~\ref{partial-periodicity-lemma}]
        Let $G := G_{k, m, q}$ and $D := D_{k, m, q}$. By Lemma~\ref{colorability-lemma} $G-v$ is $(k-1)$-colorable for every $v$, so in particular $\chi(G) \leq k$. It is left to show that every graph obtained from $G$ by removing at most $r$ edges is not $(k-1)$-colorable. Towards a contradiction, suppose that there exists a graph $G'=G-R$ with $R\subseteq E(G), |R|\le r$ that admits a proper $(k-1)$-coloring $\phi$.

        We start by showing that every set $S'$ of $\frac{N-1}{4}$ consecutive vertices has a subset of $3n_{k,m}$ consecutive unaffected vertices. Then we will use Lemma~\ref{partial-periodicity-lemma} to show that $\varphi$ is periodic. Finally, using the periodicity we will show that $\varphi$ is not proper.

        Since $G$ is symmetric with respect to cyclic shifts, we may assume without loss of generality that $S' = \{v_1, \dots, v_{\frac{N - 1}{4}}\}$ in the following. Let $a$ be the number of affected vertices in $S'$. Then $a \leq 2|R|\le 2r$. We next partition the set of unaffected vertices of $S'$ into $a + 1$ sets $U_0, \dots, U_{a}$ where $v_i \in U_j$ if and only if there are exactly $j$ affected vertices with a smaller index than $i$. Observe that each $U_j$ is a set of consecutive unaffected vertices and that their union has $\frac{N-1}{4} - a$ vertices. Then the largest set $U_i$ satisfies
        \begin{equation}\notag
            |U_i| \geq \frac{\frac{N - 1}{4} - 2r}{2r + 1} = \frac{q}{8r + 4}n_{k, m} - \frac{2r}{2r + 1} > 3n_{k, m} - 1.
        \end{equation}
        This shows that every set $S'$ of $\frac{N-1}{4}$ consecutive vertices contains a subset $S$ of $3n_{k,m}$ consecutive unaffected vertices, as desired. Applying Lemma~\ref{partial-periodicity-lemma}, we can now conclude that $\varphi$ is periodic on every set $S'$ of $\frac{N-1}{4}$ consecutive vertices.

        Now we will use the periodicity to show that $\varphi$ cannot be proper.
        From what we argued above it follows that there is some unaffected vertex in $G'$. By cyclic symmetry, we may w.l.o.g. assume that $v_0$ is unaffected. Then $v_0$ and $v_{N-1}$ are adjacent in $G$ (because $d_N(0, N-1) = 1 \in D_1$) and thus also in $G'$ (as $v_0$ is unaffected). We will now show that they are assigned the same color by $\varphi$, arriving at the desired contradiction. Since $q > 4$, we have $n_{k, m} < \frac{N - 1}{4}$. Hence, for every $i \in [1, q]$ there is a set of $\frac{N-1}{4}$ consecutive vertices containing $v_{(i-1)n_{k,m}}$ and $v_{in_{k,m}}$. Therefore, we can use the periodicity to get

        \begin{equation}\label{contradiction-eqn}
            \varphi(v_0) = \varphi(v_{n_{k,m}}) = \dots = \varphi(v_{qn_{k,m}}).
        \end{equation}

        Recall that $N-1 = qn_{k,m}$, so (\ref{contradiction-eqn}) tells us that $\varphi(v_{0}) = \varphi(v_{N-1})$. This implies that $\varphi$ is not proper, so our assumption was false. Hence, every graph obtained from $G$ by removing at most $r$ edges has chromatic number~$k$. It follows that $G_{k,m,q}$ is a $(k,r)$-graph, concluding the proof of the theorem.
    \end{proof}

\end{document}
